% Cote 119: offprints, reports, courses, CV and letters (typed and printed part): archive cote 119, site page 298 (PDF page 299).
% typescript. Draft compared: none.
% AI-transcribed by Claude Opus 5.5 (claude-opus-5-5, adaptive thinking, effort high), Message Batches API, batch r119_full (msgbatch_01DwZNm5Y7YsKBWyAN4sFs8w), 2026-11-13; not yet checked by eye.
\scanpage{119/298}
\mspage{- 11 -}

Lorsque $r$ est assez petit, ce résultat peut se traduire en termes
cohomologiques. En effet, l'hypothèse que $N$ opère transitivement sur
$S - \{1\}$ entraîne qu'il existe un élément $x$ de $H^{r}(G, Z/2Z)$ dont la
restriction à tout sous-groupe d'ordre $2$ de $G$ est $\neq 0$ , et un tel élément
est unique, à l'addition près d'une classe de cohomologie `` négli\over{g}{ge}able ''
(cf. §7 ci-après). L'élément correspondant $x_{L}$ de $H^{r}(k, Z/2Z)$ est un
invariant de l'algèbre galoisienne $L$.

(6.5.3) - \uline{Supposons} $r \leq 4$. \uline{Les conditions} (i), (ii), (iii) \uline{de} (6.5.2) \uline{sont}
\uline{alors équivalentes à} :
(iv) \uline{On a} $x_{L} = 0$ \uline{dans} $H^{r}(k, Z/2Z)$.
L'hypothèse $r \leq 4$ pourrait être supprimée si les conjectures du n°2.3
étaient démontrées.

\uline{Exemples}. 1) Supposons que $r = 2$ et que $N$ opère transitivement sur
$S - \{1\}$; c'est le cas si $G = A_{4}$ , $A_{5}$ ou $PSL_{2}(F_{q})$ avec $q \equiv 3 \pmod{8}$.
Le groupe $H^{2}(G, Z/2Z)$ contient un seul élément $x \neq 0$; soit $\overset{\hand{\sim}}{G}$ l'extension
correspondante de $G$ par $Z/2Z$. Il résulte de (6.5.3) que $L$ \uline{a une} BNA
\uline{si et seulement si l'homomorphisme} $\varphi_{L}$: $\mathrm{Gal}(k_{s}/k) \;\text{->}\; G$ \uline{se relève en un}
\uline{homomorphisme dans} $\overset{\hand{\sim}}{G}$ . Un tel relèvement correspond à une $\overset{\hand{\sim}}{G}$-algèbre
galoisienne $\overset{\hand{\sim}}{L}$; on peut montrer qu'il est possible de s'arranger pour que
$\overset{\hand{\sim}}{L}$ possède elle aussi une BNA.
2) Prenons pour $G$ le groupe $SL_{2}(F_{8})$ ou le groupe de Janko $J_{1}$ . Les
hypothèses de (6.5.2) et (6.5.3) sont alors satisfaites avec $r = 3$.
Le groupe $H^{3}(G, Z/2Z)$ contient un seul élément $x \neq 0$, et l'on voit que
$L$ \uline{a une} BNA \uline{si et seulement si} $x_{L} = 0$ \uline{dans} $H^{3}(k, Z/2Z)$.

\uline{Remarque} - La propriété pour une $G$-algèbre galoisienne $L$ d'avoir une BNA
peut se traduire en terme de `` torsion galoisienne '' de la manière suivante :
Soit $V$ un espace vectoriel de dimension finie sur $k$, muni d'une famille
$\mathbf{q} = (q_{i})$ de \uline{tenseurs quadratiques} (de type (2,0), (1,1), ou (0,2), peu
importe). Supposons que $G$ opère sur $V$ en fixant chacun des $q_{i}$. On peut alors
\uline{tordre} $(V,\mathbf{q})$ par le $G$-torseur correspondant à $L$. On obtient ainsi une
$k$-\uline{forme} $(V,\mathbf{q})_{L}$ de $(V,\mathbf{q})$. On peut démontrer :
(6.6) - \uline{Si} $L$ \uline{a une} BNA, $(V,\mathbf{q})_{L}$ \uline{est isomorphe à} $(V,\mathbf{q})$.
De plus, cette propriété \uline{caractérise} les algèbres galoisiennes ayant une
BNA.
(Noter que ce résultat serait faux pour les tenseurs cubiques.)

\mspage{\hand{298}}
